\documentclass{report}
\usepackage{amsmath}
\usepackage{graphicx}
\usepackage{hyperref}
\usepackage{listings}

\title{Internal Technical Report: Database Sharding Strategy}
\author{Infrastructure Team}
\date{Q2 2026}

\begin{document}
\maketitle

\tableofcontents

\chapter{Executive Summary}
This report evaluates three sharding strategies (range, hash, directory-based) for our primary write path.

\chapter{Background}
Our write volume has grown 4x year-over-year. Single-primary Postgres is approaching its IOPS ceiling.

\section{Current architecture}
A single Postgres 15 primary handles all writes; two replicas serve reads.

\section{Growth projections}
At current growth, we expect to exhaust the primary's capacity by late 2026.

\chapter{Evaluation}

\section{Range-based sharding}
\textbf{Pros.} Simple to reason about, easy range scans.

\textbf{Cons.} Hot shards for monotonic keys.

\section{Hash-based sharding}
\textbf{Pros.} Even load distribution.

\textbf{Cons.} Range scans require scatter-gather.

\section{Directory-based sharding}
\textbf{Pros.} Flexible shard assignment per-tenant.

\textbf{Cons.} Directory lookup overhead on every query.

\chapter{Recommendation}
We recommend hash-based sharding keyed on \texttt{tenant\_id} for the write path, with per-tenant read routing for range queries.

\chapter{Appendix}
See \url{https://wiki.internal/sharding-design} for implementation details.

\end{document}
